\documentclass[12pt,a4paper]{article}
\usepackage[utf8]{inputenc}
\usepackage[T1]{fontenc}
\usepackage[french]{babel}
\usepackage{geometry}
\usepackage{xcolor}
\usepackage{fancyhdr}

\geometry{margin=2cm}
\pagestyle{fancy}
\fancyhf{}
\fancyhead[L]{CEJM - Fiche de Révision}
\fancyhead[R]{Chapitre 5}
\fancyfoot[C]{\thepage}
\setlength{\headheight}{15pt}

% Couleurs
\definecolor{bleu}{RGB}{0,90,170}
\definecolor{vert}{RGB}{0,150,0}
\definecolor{rouge}{RGB}{200,0,0}
\definecolor{gris}{RGB}{240,240,240}

% Commandes pour encadrés
\newcommand{\definition}[2]{\noindent\fcolorbox{bleu}{gris}{\parbox{\dimexpr\textwidth-2\fboxsep-2\fboxrule\relax}{\textbf{\textcolor{bleu}{DÉFINITION :}} #1\\\textit{#2}}}}

\newcommand{\pointcle}[1]{\noindent\fcolorbox{vert}{white}{\parbox{\dimexpr\textwidth-2\fboxsep-2\fboxrule\relax}{\textbf{\textcolor{vert}{POINT CLÉ :}} #1}}}

\newcommand{\exemple}[1]{\noindent\fcolorbox{rouge}{white}{\parbox{\dimexpr\textwidth-2\fboxsep-2\fboxrule\relax}{\textbf{\textcolor{rouge}{EXEMPLE :}} #1}}}

\title{\textbf{Fiche de Révision CEJM}\\[2mm]\large Chapitre 5 : La finalité de l'entreprise et sa performance globale}
\author{}
\date{}

\begin{document}
\maketitle

\section*{Problématique}
\textbf{Comment la finalité d'une entreprise influence-t-elle sa stratégie et sa performance ?}

\section{Définitions Clés}

\definition{Finalité}{Raison d\'être de l\'entreprise, mission fondamentale qui justifie son existence sur le long terme.}

\vspace{3mm}
\definition{Responsabilité Sociétale des Entreprises (RSE)}{Intégration volontaire par les entreprises de préoccupations sociales, environnementales, et économiques dans leurs activités commerciales et leurs relations avec leurs parties prenantes.}

\vspace{3mm}
\definition{Parties prenantes (Stakeholders)}{Ensemble des acteurs internes et externes qui peuvent affecter ou être affectés par les décisions de l\'entreprise (salariés, clients, fournisseurs, actionnaires, communauté locale, etc.).}

\vspace{3mm}
\definition{Performance globale}{Résultat de l\'entreprise mesuré non seulement sur des critères économiques et financiers, but aussi sur des critères sociaux, sociétaux et environnementaux.}

\section{Les Finalités de l\'Entreprise}
\pointcle{Une entreprise ne poursuit pas uniquement un but lucratif. Ses finalités orientent sa stratégie.}

\subsection{La finalité économique}
\begin{itemize}
    \item Produire des biens ou des services pour un marché.
    \item Assurer sa pérennité et sa croissance (rentabilité, profit).
    \item Créer de la valeur et des richesses.
\end{itemize}

\subsection{La finalité sociale}
\begin{itemize}
    \item Concerne la gestion des ressources humaines et les conditions de travail.
    \item \textbf{Exemples :} politique de rémunération, formation, sécurité, bien-être au travail, insertion de personnes en difficulté.
\end{itemize}

\subsection{La finalité environnementale}
\begin{itemize}
    \item Vise à limiter l\'impact de l\'activité de l\'entreprise sur l\'environnement.
    \item \textbf{Exemples :} réduction des émissions de CO2, utilisation de matériaux recyclés/durables, économie d\'énergie, gestion des déchets.
\end{itemize}

\subsection{La finalité sociétale}
\begin{itemize}
    \item Concerne l\'impact de l\'entreprise sur la société au sens large.
    \item \textbf{Exemples :} soutien aux communautés locales, éthique des affaires, respect des droits de l\'homme, contribution à des causes d\'intérêt général.
\end{itemize}

\exemple{L\'entreprise \textbf{Véja} illustre ces quatre finalités :
\begin{itemize}
    \item \textbf{Économique :} Vente de baskets pour être rentable.
    \item \textbf{Sociale :} Bonnes conditions de travail et insertion de personnes en situation de handicap.
    \item \textbf{Environnementale :} Utilisation de coton agroécologique et de caoutchouc d\'Amazonie.
    \item \textbf{Sociétale :} Juste rémunération des producteurs et préfinancement des récoltes.
\end{itemize}}

\section{La Démarche RSE et les Parties Prenantes}
\pointcle{La RSE est la mise en œuvre concrète des finalités sociales, environnementales et sociétales.}

\subsection{Les bénéfices d\'une démarche RSE}
\begin{itemize}
    \item \textbf{Image de marque :} Améliore la réputation et l\'attractivité de l\'entreprise.
    \item \textbf{Fidélisation :} Attire et retient les talents (marque employeur) et les clients.
    \item \textbf{Innovation :} Pousse à trouver de nouvelles solutions plus durables.
    \item \textbf{Gestion des risques :} Anticipe les risques réglementaires, sociaux et environnementaux.
    \item \textbf{Performance économique :} Peut entraîner des économies (énergie, matières premières) et ouvrir de nouveaux marchés.
\end{itemize}

\subsection{Répondre aux attentes des parties prenantes}
\pointcle{Une stratégie RSE efficace identifie les attentes des parties prenantes et y répond de manière équilibrée.}
\begin{itemize}
    \item \textbf{Clients :} Attendent des produits de qualité, éthiques et transparents.
    \item \textbf{Salariés :} Recherchent du sens, de bonnes conditions de travail et une juste rémunération.
    \item \textbf{Fournisseurs :} Veulent des relations commerciales équitables et durables.
    \item \textbf{Actionnaires :} Sont de plus en plus attentifs à la performance globale et aux risques extra-financiers.
    \item \textbf{Société civile (ONG, communautés)} : Exige que l\'entreprise assume sa responsabilité et contribue positivement à la société.
\end{itemize}

\section{La Performance Globale}
\pointcle{La performance ne se limite pas au profit. C\'est la capacité à concilier les dimensions économiques, sociales et environnementales.}
\begin{itemize}
    \item \textbf{Indicateurs économiques :} Chiffre d\'affaires, résultat net, rentabilité (voir Chapitre 4).
    \item \textbf{Indicateurs sociaux :} Taux de turn-over, budget formation, indice d\'égalité hommes-femmes.
    \item \textbf{Indicateurs environnementaux :} Bilan carbone, part des énergies renouvelables, pourcentage de déchets recyclés.
    \item \textbf{Indicateurs sociétaux :} Montant des achats au secteur solidaire, mécénat.
\end{itemize}

\exemple{La performance de \textbf{Véja} est globale car l\'entreprise est rentable tout en ayant un impact social et environnemental positif, ce qui répond aux attentes fortes de ses clients et partenaires.}

\section{Mots-clés à connaître}
\begin{itemize}
    \item Finalité (économique, sociale, environnementale, sociétale)
    \item Responsabilité Sociétale des Entreprises (RSE)
    \item Parties prenantes (stakeholders)
    \item Performance globale
    \item Éthique des affaires
    \item Développement durable
    \item Marque employeur
\end{itemize}

\end{document}
